\section{The Congressional--State Legislative Distinction}
\label{sec:congressional}

\subsection{Article I and Congressional Redistricting}

The constitutional basis for congressional redistricting differs materially from state legislative redistricting, and this difference has direct bearing on the permissibility of citizen-VAP redistricting at the congressional level.

Article I, Section 2 of the Constitution provides that Representatives shall be ``apportioned among the several States according to their respective Numbers''---a phrase interpreted throughout American history to mean total population, not citizen population. The Fourteenth Amendment's Apportionment Clause (Section 2) reinforced this reading by specifying the apportionment base as ``the whole number of persons in each State,'' while separately penalizing states that deny the vote to male citizens---a distinction that would be meaningless if the apportionment base were limited to citizens.

This textual history supports the view that congressional apportionment---and by extension, congressional redistricting, which must achieve equal population \emph{within} a state's apportionment---is constitutionally anchored to total population. Under this view:
\begin{enumerate}
    \item Federal apportionment is based on total population (uncontested since the founding)
    \item A state's congressional seats are allocated based on total population
    \item Congressional districts within the state must be approximately equal in total population to comply with \emph{Wesberry v. Sanders}~\citep{wesberry1964}
    \item Using citizen VAP as the intrastate balance constraint would produce districts of unequal total population---which Congress could potentially challenge under its Article I authority over federal elections
\end{enumerate}

No Supreme Court decision has directly held that states \emph{may not} use citizen VAP for congressional redistricting. But the textual and structural argument that Article I requires total-population congressional districts is stronger than the parallel argument for state legislative redistricting, where the constitutional hook is the more open-ended Equal Protection Clause rather than the specific Apportionment Clause text.

\subsection{State Legislative Redistricting After Evenwel}

For state legislative redistricting, \emph{Evenwel} controls. The decision's core holding---that states \emph{may} use total population---is narrow: it answers one direction of the question (challengers cannot \emph{force} citizen-VAP), while leaving open the other (challengers might \emph{choose} citizen-VAP). The two open questions are:

\begin{enumerate}
    \item \textbf{CVAP-choice permissibility}: If a state legislature elects to use citizen VAP for state legislative redistricting, does that violate the Equal Protection Clause by systematically disadvantaging noncitizen-dense communities? The lower courts are divided; no Supreme Court resolution.

    \item \textbf{Discriminatory purpose}: If a state adopts citizen-VAP redistricting with demonstrable intent to disadvantage Hispanic communities (which have the largest noncitizen fractions), does that constitute intentional discrimination under \emph{Washington v. Davis}~\citep{washingtondavis1976} or \emph{Village of Arlington Heights v. Metropolitan Housing Development Corp.}~\citep{arlingtonheights1977}? This question is separate from discriminatory effect analysis.
\end{enumerate}

\subsection{Burns v. Richardson and the CVAP Analogy}

\emph{Burns v. Richardson}, 384 U.S. 73 (1966)~\citep{burns1966} is the closest precedent for state use of a non-total-population base. Hawaii used registered voters to draw its state senate districts, excluding large numbers of military personnel counted in the census but not registered in Hawaii. The Court approved this approach, emphasizing that registered voters were not a \emph{systematically discriminatory} exclusion---the military population was transient and had minimal attachment to Hawaii.

The \emph{Burns} rationale does not straightforwardly extend to citizen-VAP redistricting because noncitizens are not transient and do have local attachments. Unlike military personnel (who maintain legal domicile elsewhere), noncitizens may have resided in the United States for decades, have U.S.-born citizen children, own property, and participate in local economic and civic life. The \emph{Burns} Court's emphasis on transience as the justification for exclusion suggests that a citizen-VAP base, which excludes long-term residents, presents a harder constitutional question than \emph{Burns} addressed.

\subsection{The Deliberate Exclusion Problem}

The deepest constitutional problem with mandatory citizen-VAP redistricting is not the mechanical effect on district boundaries but the deliberate exclusion of permanent residents from political representation. If noncitizens are excluded from the population base, their representative will be allocated to represent fewer actual constituents---not because the representative will ignore them, but because the district will be drawn to dilute their effective political influence by combining them with more-citizen suburban communities.

This is the representation-theory argument that the \emph{Evenwel} majority found compelling in the opposite direction: the Court held that noncitizens \emph{are} legitimate constituents whose representation interests justify their inclusion in the population base. A future decision holding that states \emph{must} exclude noncitizens would need to explain why noncitizens' representation interests---which the \emph{Evenwel} Court recognized---do not count. That explanation is not available in existing Equal Protection doctrine.
